\documentclass[12pt]{article}
\usepackage{amsmath, amssymb}
\usepackage{geometry}
\geometry{margin=1in}

\title{Boat-River Vector Simulation Using Custom Vector Engine}
\author{Alomgir Bhuiyan}
\date{\today}

\begin{document}

\maketitle

\section{Introduction}

This document explains a C-based vector simulation program that models boat motion in a river using a custom vector library. The system uses both Cartesian and Polar vector representations to compute resultant velocity and direction.

The goal is to determine:
\begin{itemize}
    \item Final velocity of the boat
    \item Direction (angle) of motion relative to the river bank
\end{itemize}

\section{Physical Model}

The motion is based on relative velocity:

\[
\vec{R} = \vec{V}_{boat} + \vec{V}_{river}
\]

Where:
\begin{itemize}
    \item $\vec{V}_{boat}$ = velocity of boat in still water
    \item $\vec{V}_{river}$ = velocity of river current
    \item $\vec{R}$ = resultant velocity
\end{itemize}

\section{Vector Engine Usage}

The program uses a custom vector library with:
\begin{itemize}
    \item Cartesian vectors: $(i, j, k)$
    \item Polar vectors: $(r, \theta)$ or $(r, \theta_x, \theta_y, \theta_z)$
    \item Workspace memory system for allocation
\end{itemize}

\section{Memory Allocation System}

\[
\text{VectorWorkspace} = \text{preallocated memory block}
\]

Instead of multiple allocations, all vectors are stored inside a single workspace:

\[
\texttt{VectorAllocate(24)}
\]

This improves performance and reduces fragmentation.

\section{Boat and River Velocity Definition}

The velocities are defined in polar form:

\subsection{Boat Velocity}

\[
V_{boat} = (r = 4, \theta = 120^\circ)
\]

This means the boat moves with magnitude 4 at an angle of 120 degrees.

\subsection{River Velocity}

\[
V_{river} = (r = 2, \theta = 0^\circ)
\]

This represents current flow along a fixed axis.

\section{Conversion to Cartesian Form}

The library converts polar vectors to Cartesian form using:

\[
V_x = r \cos\theta
\]
\[
V_y = r \sin\theta
\]

So:
\begin{itemize}
    \item Boat vector becomes $(x_b, y_b)$
    \item River vector becomes $(x_r, y_r)$
\end{itemize}

\section{Resultant Vector Calculation}

The resultant vector is computed using vector addition:

\[
\vec{R} = \vec{V}_{boat} + \vec{V}_{river}
\]

Component-wise:

\[
R_x = V_{bx} + V_{rx}
\]
\[
R_y = V_{by} + V_{ry}
\]

This is implemented using:

\texttt{Resultant(vectors, 2, vws)}

\section{Conversion Back to Polar Form}

After computing the resultant vector in Cartesian form, it is converted back:

\[
r = \sqrt{R_x^2 + R_y^2}
\]

\[
\theta = \tan^{-1}\left(\frac{R_y}{R_x}\right)
\]

This is handled by:

\texttt{CartesianVectorToPolar(R, vws)}

\section{Final Output Interpretation}

The program prints:

\begin{itemize}
    \item Magnitude of resultant velocity
    \item Angle of motion relative to river bank
\end{itemize}

\[
\text{Final velocity} = |\vec{R}|
\]
\[
\text{Direction} = \theta
\]

\section{Why This Design Works}

This system is designed to:
\begin{itemize}
    \item Separate physics (math) from implementation (C code)
    \item Support both coordinate systems
    \item Allow reusable vector operations
    \item Model real-world motion using vector algebra
\end{itemize}

\section{Conclusion}

The program successfully simulates boat-river motion using a custom vector engine. It demonstrates real-world application of vector addition, coordinate conversion, and relative velocity using a structured computational approach.

\end{document}
